\chapter{较大素数的奇次幂：已经做出了什么，还缺什么}\label{ch:high}
\section{旧周期消失，不等于问题无解}
现在令 $k=ps^2,r=ps,n=k+r+1$，其中 $p\ge5$。若要研究 $p^{2a+1}$，就取 $s=p^a$。上一讲的基本返回有 $\theta=p>4$，没有有限射影周期。

继续研究的方向是换返回词。当前已找到两类三步返回，其整数矩阵迹为零，因此每个返回重复两次都是标量。它们给出了通用的正向可逆资源，但尚未产生完整剪切群。下面保留其构造，因为它准确展示了进展和剩余的容量问题。\src{S20}

显示的计数与矩阵恒等式实际上允许任意整数 $p\ge5$；本讲关心的应用是 $p$ 为素数的情形。

\section{用三次平均制造一个反射}
假设
\[s\ge(p-1)(p-2).
\]
从核心 $K=(u^k,v^r,w)$ 出发，先取 $j=s-1$，平均
\[(u^{k-j-1},v^j,w).
\]
得到 $k$ 份 $D$，外面剩 $s$ 份旧 $u$ 和 $d=(p-1)s+1$ 份旧 $v$。第二步取 $a$ 份旧 $u$、$b$ 份旧 $v$ 和 $k-a-b$ 份 $D$，得到 $k$ 份 $E$。

第三步保留 $r$ 份 $E$ 和一个指定旧单点，平均其余 $k$ 项。两种可用的第二步计数为
\[
\begin{array}{c|cc}
\text{最终保留}&a&b\\\hline
u&s-(p-1)(p-3)&(p-1)s+3-p\\
v&s-(p-1)(p-2)&(p-1)s+2-p
\end{array}
\]
条件 $s\ge(p-1)(p-2)$ 确保这些数非负；同时它们不超过第一步剩下的库存，而且仍留有所要求的旧单点。这些不等式正是三步词能在原位置上实现的依据。

\section{矩阵公式与它真正证明的范围}
第一行计数对应 $h=p(s-p+3)$，第二行对应 $h=p(s-p+2)$。令
\[\alpha=nh-k(r+1),\qquad\beta=k(d-b)-dh.
\]
三步作用为 $J/k^3$，其中
\[
J=\begin{pmatrix}
-r\alpha-k^2&-r\beta+z\\
-(r+k)\alpha-k^2&-(r+k)\beta+z
\end{pmatrix},
\]
保留 $u$ 时 $z=0$，保留 $v$ 时 $z=k^2$。代入两行计数都得到 $\operatorname{tr}J=0$，所以
\[J^2=\lambda I,\qquad \lambda\equiv k^6\pmod n.
\]
由于 $\gcd(k,n)=1$，$\lambda$ 非零且模 $n$ 可逆。于是两类三步返回各自有正向射影逆。这是对整族参数的恒等式证明。

对真正的奇素数幂，若 $s=p^a$ 且 $a\ge2$，容量条件成立，所以覆盖指数 $2a+1\ge5$。但立方对应 $s=p$，对 $p\ge5$ 不满足这个条件。比如 $p=5$ 要求 $s\ge12$，而立方只有 $s=5$。

\section{两块尚未补齐的内容}
第一，三次幂需要新的计数方案，或允许改变核心形状。现有表中出现负数就意味着没有那么多可选的位置，不能把它视作形式矩阵上的小问题。

第二，即使已经有两个可逆反射，也还要证明它们和其他返回能产生足够的上、下剪切与倍率。两个反射可以生成很小的群；可逆性本身不保证同余群覆盖，更不保证每个合法列都能到终端。

因此当前准确状态是：高奇次幂有两种统一的可逆三步资源；完整核心控制、全输入收尾和精确临界尚未完成。前面的GRH小单位定理不能补这一步，因为它只能判断\emph{已经实际供应的}倍率是否足够。
